% simplicial_complex_embedders_v2.tex
%
% A rewrite of v1 that presents intent first and algorithm second.
% The algorithm steps are introduced as the natural implementation of
% a geometric strategy, not as a sequence of operations to be explained
% after the fact.

\documentclass[11pt,a4paper]{article}

\usepackage[margin=2.5cm]{geometry}
\usepackage{amsmath,amssymb}
\usepackage{booktabs}
\usepackage{array}
\usepackage{xcolor}
\usepackage{colortbl}
\usepackage{microtype}
\usepackage{enumitem}
\usepackage{hyperref}
\usepackage[most]{tcolorbox}

% ── Colours ──────────────────────────────────────────────────────────────────
\definecolor{alg1col}{HTML}{1a5276}
\definecolor{alg2col}{HTML}{1e8449}
\definecolor{boxbg}{HTML}{f4f6f7}
\definecolor{notebg}{HTML}{fef9e7}
\definecolor{keybg}{HTML}{eaf4fb}
\definecolor{warnbg}{HTML}{fdfefe}

\tcbset{
  defbox/.style ={colback=boxbg,  colframe=gray!50,   boxrule=0.4pt,
                  arc=2pt, left=6pt, right=6pt, top=4pt, bottom=4pt},
  notebox/.style={colback=notebg, colframe=orange!60, boxrule=0.4pt,
                  arc=2pt, left=6pt, right=6pt, top=4pt, bottom=4pt},
  keybox/.style ={colback=keybg,  colframe=blue!45,   boxrule=0.7pt,
                  arc=3pt, left=8pt, right=8pt, top=6pt, bottom=6pt},
}

% ── Notation ─────────────────────────────────────────────────────────────────
\newcommand{\eji}[2]{e^{#1}_{#2}}       % Eji: vector index raised, coord lowered
\newcommand{\hash}[1]{h(#1)}            % hash to hypercube
\newcommand{\canon}[1]{\widetilde{#1}}  % canonical form
\newcommand{\R}{\mathbb{R}}
\newcommand{\SP}[2]{SP^{#1}(\R^{#2})}  % symmetric product

\hypersetup{colorlinks=true, linkcolor=alg1col, urlcolor=alg1col}

% ─────────────────────────────────────────────────────────────────────────────
\title{\textbf{Simplicial Complex Multiset Embedders}\\[0.4em]
       \large Algorithms 1 and 2 — Strategy and Implementation (v2)}
\date{\today}
% ─────────────────────────────────────────────────────────────────────────────

\begin{document}
\maketitle
\tableofcontents
\bigskip

% ═════════════════════════════════════════════════════════════════════════════
\section{What We Are Trying to Do}
% ═════════════════════════════════════════════════════════════════════════════

The input is an unordered multiset $M = \{\mathbf{v}_0,\ldots,\mathbf{v}_{n-1}\}$
of $n$ vectors in $\R^k$.  Concretely we represent $M$ as a $k\times n$ matrix
whose columns are the vectors (column $j$ = vector $\mathbf{v}_j$, row $i$ =
coordinate $i$).  Column order is arbitrary.

We want a map $f$ from such multisets to some $\R^N$ that is:
\begin{itemize}[nosep]
  \item \textbf{Multiset-invariant} — the output does not depend on how the
        columns are ordered.
  \item \textbf{Injective} — distinct multisets give distinct outputs ($f$ is
        an \emph{embedder}).
  \item \textbf{Continuous} ($C^0$) — small changes in the vectors produce
        small changes in $f(M)$.
  \item \textbf{Positively homogeneous of degree 1} — $f(\lambda M) =
        \lambda f(M)$ for all $\lambda > 0$.
\end{itemize}

Multiset invariance and continuity are in tension: any map that achieves
invariance by sorting will have kinks wherever the sorted order changes.
The strategy below resolves this tension elegantly.

% ═════════════════════════════════════════════════════════════════════════════
\section{The Strategy}
\label{sec:strategy}
% ═════════════════════════════════════════════════════════════════════════════

\subsection*{Step A — Write $M$ in a natural basis}

Label each of the $nk$ matrix entries by $\eji{j}{i}$ (column $j$, row $i$).
Let $X[j][i]$ denote the entry in column $j$, row $i$.  Then
\[
  M \;=\; \sum_{j,i} X[j][i]\;\eji{j}{i},
\]
where $\eji{j}{i}$ is the $k\times n$ matrix that is 1 at $(i,j)$ and 0
elsewhere.  This is the expansion of $M$ in the \emph{standard Eji basis}.

\subsection*{Step B — Peel off the fully symmetric parts}

A \emph{fully symmetric} basis element at row $i$ is $\sum_j \eji{j}{i}$: it
places the same coefficient on every vector at coordinate $i$, so it is
unaffected by any permutation of the columns.  The smallest coefficient that
any vector has at row $i$ is $m_i = \min_j X[j][i]$.  Factoring:
\[
  M \;=\; \underbrace{\sum_i m_i \!\left(\sum_j \eji{j}{i}\right)}_{\text{symmetric part}}
    \;+\; \underbrace{\text{balance}}_{\text{all entries} \;\ge\; 0}.
\]
The symmetric part is already permutation-invariant; its $k$ scalar
coefficients $m_0,\ldots,m_{k-1}$ can be stored directly as anchor values.

The balance $M - \sum_i m_i (\sum_j \eji{j}{i})$ has all non-negative
entries and contains the non-trivial, non-symmetric information in $M$.

\subsection*{Step C — Re-express the balance as a positive combination of
             nested basis elements}

Within row $i$, order the $n$ vectors descending by their (shifted) value.
Let $\sigma_i(0),\ldots,\sigma_i(n-1)$ be this ordering.  The gaps
$\delta^{(i)}_r = X[\sigma_i(r)][i] - X[\sigma_i(r+1)][i] \ge 0$ are the
drops between consecutive ranks.

By Abel summation (the discrete analogue of integration by parts), the balance
at row $i$ becomes:
\[
  \text{balance at row }i
    \;=\; \sum_{r=0}^{n-2} \delta^{(i)}_r \cdot V^{(i)}_r,
\]
where the \emph{prefix basis element} $V^{(i)}_r =
\bigl\{\eji{\sigma_i(0)}{i},\ldots, \eji{\sigma_i(r)}{i}\bigr\}$
is the set of the top-$(r{+}1)$ vectors by row-$i$ value.

This is another change of basis: from individual Ejis with signed coefficients
to nested prefix sets with non-negative gap coefficients.  The structure of
which vectors appear in $V^{(i)}_r$ — which $j$'s dominate row $i$ — is the
non-symmetric information not yet discarded.

\subsection*{Step D — Merge, re-sort, and apply a second Abel summation}

Collect all $m = k(n-1)$ gap--prefix-element pairs from all rows into one list
and sort descending by gap: $\delta_{(0)} \ge \cdots \ge \delta_{(m-1)}$ with
corresponding basis elements $V_{(0)},\ldots,V_{(m-1)}$.

Apply a second Abel summation to obtain \emph{cumulative} basis elements and
weighted-difference coefficients:
\[
  \alpha_s \;=\; (s+1)\bigl(\delta_{(s)} - \delta_{(s+1)}\bigr),
  \qquad \delta_{(m)} := 0,
\]
\[
  L_s \;=\; V_{(0)} + V_{(1)} + \cdots + V_{(s)},
\]
where $L_s$ is the \emph{Eji\_LinComb}: the $k\times n$ matrix whose $(i,j)$
entry counts how many of the first $s+1$ prefix elements contain $\eji{j}{i}$.
Its \emph{index} $s+1$ records how many prefix elements have been accumulated.

The mathematical identity $\sum_s \delta_{(s)} V_{(s)} = \sum_s \alpha_s
\cdot (L_s/(s+1))$ confirms this preserves all information.  The normalised
basis element $L_s/(s+1)$ is the \emph{average} of the first $s+1$ prefix
elements.  The $(s+1)$ weight in $\alpha_s$ distributes contributions more
evenly and ensures that $\sum_s \alpha_s = \sum_s \delta_{(s)}$ (the total
weight is preserved).

Crucially, $L_s$ captures \emph{which prefix elements came first} in the
global sort: two different orderings of the same set of prefix elements produce
different $L_s$ sequences.

\subsection*{Step E — Achieve permutation invariance by canonicalising}

The $j$-indices in $L_s$ still reflect the arbitrary labelling of columns.
The \emph{canonical form} $\canon{L_s}$ is obtained by sorting the columns of
$L_s$'s count matrix lexicographically — this removes the column labelling
while retaining the multiset structure.  Two inputs related by a column
permutation produce identical $\canon{L_s}$ sequences and hence identical
outputs.

\subsection*{Step F — Encode in $\R^{2m+1}$ and assemble}

Map each canonical $L_s$ to a pseudorandom point $p_s = \hash{\canon{L_s}}
\in [0,1]^{2m+1}$ via MD5 hashing of its count matrix and index.  The main
body of the output is $\sum_s \alpha_s p_s$.

\begin{tcolorbox}[keybox]
\textbf{Why $\R^{2m+1}$ is the right dimension.}
All $\alpha_s \ge 0$, so $\sum_s \alpha_s p_s$ is a point in the
\emph{positive cone} generated by $\{p_0,\ldots,p_{m-1}\}$, a cone of
dimension $m$.  Two $m$-dimensional positive cones in $\R^{2m+1}$ are
\emph{generically disjoint} by dimension counting: $m + m = 2m < 2m+1$.
Therefore a single point in $\R^{2m+1}$ encodes both:
(a) which cone it lies in (i.e.\ which canonical $\{L_s\}$, i.e.\ which basis
elements), and (b) the coefficients $\{\alpha_s\}$ within that cone.

This is also why the algorithm is $C^0$: when two values in the same row are
equal, the gap at that position is $0$ and so $\alpha_s = 0$.  The point
$\sum \alpha_s p_s$ then lies on a shared \emph{face} of two adjacent cones
rather than in either interior.  That shared face is the geometric expression
of the continuity: the change in which cone we are in happens exactly at the
moment when the coefficient of the transitioning basis element is zero, so the
output is unaffected.
\end{tcolorbox}

\noindent
The full output is:
\[
  f(M) \;=\; \Bigl[\;\underbrace{m_0,\ldots,m_{k-1}}_{k\;\text{row minima}},\;
                     \underbrace{\textstyle\sum_{s=0}^{m-1}\alpha_s p_s}_{2m+1\;\text{reals}}
             \;\Bigr]
  \;\in\; \R^{k + 2m + 1}.
\]
With $m = k(n-1)$: total size $= k + 2k(n-1) + 1 = 2nk+1-k$.

% ═════════════════════════════════════════════════════════════════════════════
\section{Two Algorithms, One Strategy}
\label{sec:two-algs}
% ═════════════════════════════════════════════════════════════════════════════

Both algorithms follow Steps A--F.  They differ only in \textbf{Step B}: how
many symmetric parts are peeled off, and what ``fully symmetric'' means.

\paragraph{Algorithm 2 (per-row).}
Factors out the $k$ per-row symmetric elements $\sum_j \eji{j}{i}$, one per
coordinate.  Stores $k$ row minima as anchors.  Steps C--F operate on
$m = k(n-1)$ gaps.

\paragraph{Algorithm 1 (global).}
Factors out only the \emph{single} fully-symmetric element $\sum_{i,j}
\eji{j}{i}$ (treating every entry identically), with coefficient = global
minimum.  The global maximum is also recorded (see note below).  Steps C--F
then operate on a single globally-sorted list of $m = nk-1$ gaps.

The single global sort of Algorithm~1 means that prefix elements can contain
Ejis from \emph{any} row, whereas in Algorithm~2 every prefix element contains
Ejis from exactly one row.  This is why Algorithm~2 is conceptually cleaner:
it respects the coordinate structure of the vectors.

\paragraph{Dimension count.}
\begin{center}
\renewcommand{\arraystretch}{1.3}
\begin{tabular}{lcc}
  \toprule
  & \textbf{\textcolor{alg1col}{Algorithm 1}} & \textbf{\textcolor{alg2col}{Algorithm 2}} \\
  \midrule
  Symmetric parts extracted & 1 (global) & $k$ (per-row) \\
  Gaps $m$ & $nk-1$ & $k(n-1)=nk-k$ \\
  Anchor values stored & 2 (max \& min) & $k$ (row minima) \\
  Hash embedding dimension $2m+1$ & $2nk-1$ & $2nk-2k+1$ \\
  \textbf{Total output size} & $\mathbf{2nk+1}$ & $\mathbf{2nk+1-k}$ \\
  \bottomrule
\end{tabular}
\end{center}

Algorithm~2 saves $k$ output dimensions because per-row symmetric extraction
produces $k-1$ fewer gaps than global extraction (since each row contributes
$n-1$ gaps, totalling $k(n-1) = nk-k$, versus the global $nk-1$ gaps from
Algorithm~1), and the $(s+1)$ weighting scheme in Step~D ensures that the
hash dimension savings exceed the extra anchor overhead.

\begin{tcolorbox}[notebox]
\textbf{Note on Algorithm 1's global maximum.}  In a perfect implementation,
only the global minimum is needed as an anchor (it fixes the absolute scale;
all relative gaps are encoded in the hash sum).  The global maximum is
technically redundant.  It is included in Algorithm~1 for implementation
simplicity; a TODO comment in the source code acknowledges this.
\end{tcolorbox}

% ═════════════════════════════════════════════════════════════════════════════
\section{Worked Example: Algorithm 2, $n=4$, $k=3$}
\label{sec:example-alg2}
% ═════════════════════════════════════════════════════════════════════════════

\[
  M \;=\; \left\{\;
    \begin{pmatrix}4\\2\\3\end{pmatrix},\;
    \begin{pmatrix}-3\\5\\1\end{pmatrix},\;
    \begin{pmatrix}8\\9\\2\end{pmatrix},\;
    \begin{pmatrix}2\\7\\2\end{pmatrix}
  \;\right\}
  \quad (n=4 \text{ vectors in } \R^3)
\]

\subsection*{Step B: Extract row minima}

Row minima: $m_0 = -3$ (row~0), $m_1 = 2$ (row~1), $m_2 = 1$ (row~2).
These three scalars are stored as anchors.  After subtracting:
\[
  \text{balance} = \left\{\;
    \begin{pmatrix}7\\0\\2\end{pmatrix},\;
    \begin{pmatrix}0\\3\\0\end{pmatrix},\;
    \begin{pmatrix}11\\7\\1\end{pmatrix},\;
    \begin{pmatrix}5\\5\\1\end{pmatrix}
  \;\right\}
\]
(All non-negative.)  The debug check in the source code verifies that
$\sum_s \alpha_s \cdot (L_s/(s+1)) + \text{broadcast}(m_i)$ recovers
the original $M$ exactly.

\subsection*{Step C: Per-row Abel summation}

\noindent
For each row, sort vectors descending and compute gaps:

\begin{center}
\small
\renewcommand{\arraystretch}{1.2}
\begin{tabular}{clll}
\toprule
Row $i$ & Sorted descent (value: vector $j$) & Gaps $\delta$ & Prefix elements $V^{(i)}_r$ \\
\midrule
0 & $11(j{=}2),\;7(j{=}0),\;5(j{=}3),\;0(j{=}1)$ & 4,\;2,\;5 &
  $\{\eji{2}{0}\}$;\;$\{\eji{2}{0},\eji{0}{0}\}$;\;$\{\eji{2}{0},\eji{0}{0},\eji{3}{0}\}$ \\[2pt]
1 & $7(j{=}2),\;5(j{=}3),\;3(j{=}1),\;0(j{=}0)$ & 2,\;2,\;3 &
  $\{\eji{2}{1}\}$;\;$\{\eji{2}{1},\eji{3}{1}\}$;\;$\{\eji{2}{1},\eji{3}{1},\eji{1}{1}\}$ \\[2pt]
2 & $2(j{=}0),\;1(j{=}2),\;1(j{=}3),\;0(j{=}1)$ & 1,\;0,\;1 &
  $\{\eji{0}{2}\}$;\;$\{\eji{0}{2},\eji{2}{2}\}$;\;$\{\eji{0}{2},\eji{2}{2},\eji{3}{2}\}$ \\
\bottomrule
\end{tabular}
\end{center}

\noindent
At this stage in the code, each pair is stored as $(\delta,\,\texttt{MSV})$
where the MSV is an individual prefix set — it is a valid basis element but
not yet in its final cumulative form.

\subsection*{Steps D: Merge, sort, and second Abel summation}

Merge all $9 = k(n{-}1)$ pairs and sort by gap (i.e.~$\delta$) descending:
\[
  5,\;4,\;3,\;2,\;2,\;2,\;1,\;1,\;0
\]
The table below shows to which basis element
$V_{(s)}$ each sorted gap $\delta_{(s)}$ was associated, and
accumulates these $V_{(s)}$ into $L_s = V_{(0)}+\cdots+V_{(s)}$.
Terms within each $L_s$ expression are grouped by coordinate row $i$.

\begin{center}
\small
\setlength{\tabcolsep}{4pt}
\renewcommand{\arraystretch}{1.35}
\begin{tabular}{c|cp{3cm}|cp{3cm}|cp{3cm}}
\toprule
    &
    &
    & $\delta_{(s)}-\delta_{(s+1)}$
    & $L_{(s-1)}+V_{(s)}$
    & $(s+1)\Delta_{(s)}$
    & $L_{(s)}/(s+1)$\\
\midrule
$s$ & $\delta_{(s)}$
    & $V_{(s)}$
    & $\Delta_{(s)}$
    & $L_{(s)}$
    & $\alpha_s$
    & ${\hat L}_{(s)}$ \\
\midrule
0 & 5 & $\eji{0}{0}+\eji{2}{0}+\eji{3}{0}$ & 1
  & $\eji{0}{0}+\eji{2}{0}+\eji{3}{0}$
  & 1
  & $\eji{0}{0}+\eji{2}{0}+\eji{3}{0}$
  \\[3pt]
1 & 4 & $\eji{2}{0}$ & 1
  & $\eji{0}{0}+2\eji{2}{0}+\eji{3}{0}$
  & 2
  & $\tfrac 1 2 (\eji{0}{0}+2\eji{2}{0}+\eji{3}{0})$
 \\[3pt]
2 & 3 & $\eji{1}{1}+\eji{2}{1}+\eji{3}{1}$
& 1
  & $\eji{0}{0}+2\eji{2}{0}+\eji{3}{0}$\newline
    ${}+\eji{1}{1}+\eji{2}{1}+\eji{3}{1}$
    & 3
  & $\tfrac 1 3 (\eji{0}{0}+2\eji{2}{0}+\eji{3}{0}$\newline
    ${}+\eji{1}{1}+\eji{2}{1}+\eji{3}{1})$
   \\[3pt]
\rowcolor{gray!15}\cellcolor{white}3 & \textcolor{gray}{$2$} & \textcolor{gray}{$\eji{0}{0}+\eji{2}{0}$} & \textcolor{gray}{$0$}
  & \textcolor{gray}{$2\eji{0}{0}+3\eji{2}{0}+\eji{3}{0}$\newline
    ${}+\eji{1}{1}+\eji{2}{1}+\eji{3}{1}$}
    & \textcolor{gray}{$0$}
  & \textcolor{gray}{$\tfrac 1 4 (2\eji{0}{0}+3\eji{2}{0}+\eji{3}{0}$\newline
    ${}+\eji{1}{1}+\eji{2}{1}+\eji{3}{1})$}
   \\[3pt]
\rowcolor{gray!15}\cellcolor{white}4 & \textcolor{gray}{$2$} & \textcolor{gray}{$\eji{2}{1}$} & \textcolor{gray}{$0$}
  & \textcolor{gray}{$2\eji{0}{0}+3\eji{2}{0}+\eji{3}{0}$\newline
    ${}+\eji{1}{1}+2\eji{2}{1}+\eji{3}{1}$}
    & \textcolor{gray}{$0$}
  & \textcolor{gray}{$\tfrac 1 5 (2\eji{0}{0}+3\eji{2}{0}+\eji{3}{0}$\newline
    ${}+\eji{1}{1}+2\eji{2}{1}+\eji{3}{1})$}
    \\[3pt]
5 & \cellcolor{gray!15}\textcolor{gray}{$2$} & \cellcolor{gray!15}\textcolor{gray}{$\eji{2}{1}+\eji{3}{1}$} & 1
  & $2\eji{0}{0}+3\eji{2}{0}+\eji{3}{0}$\newline
    ${}+\eji{1}{1}+3\eji{2}{1}+2\eji{3}{1}$
    & 6
  & $\tfrac 1 6 (2\eji{0}{0}+3\eji{2}{0}+\eji{3}{0}$\newline
    ${}+\eji{1}{1}+3\eji{2}{1}+2\eji{3}{1})$
   \\[3pt]
\rowcolor{gray!15}\cellcolor{white}6 & \textcolor{gray}{$1$} & \textcolor{gray}{$\eji{0}{2}$} & \textcolor{gray}{$0$}
  & \textcolor{gray}{$2\eji{0}{0}+3\eji{2}{0}+\eji{3}{0}$\newline
    ${}+\eji{1}{1}+3\eji{2}{1}+2\eji{3}{1}$\newline
    ${}+\eji{0}{2}$}
    & \textcolor{gray}{$0$}
 & \textcolor{gray}{$\tfrac 1 7 (2\eji{0}{0}+3\eji{2}{0}+\eji{3}{0}$\newline
    ${}+\eji{1}{1}+3\eji{2}{1}+2\eji{3}{1}$\newline
    ${}+\eji{0}{2})$}
     \\[3pt]
7 & \cellcolor{gray!15}\textcolor{gray}{$1$} & \cellcolor{gray!15}\textcolor{gray}{$\eji{0}{2}+\eji{2}{2}+\eji{3}{2}$} & 1
  & $2\eji{0}{0}+3\eji{2}{0}+\eji{3}{0}$\newline
    ${}+\eji{1}{1}+3\eji{2}{1}+2\eji{3}{1}$\newline
    ${}+2\eji{0}{2}+\eji{2}{2}+\eji{3}{2}$
    & 8
  & $\tfrac 1 8 (2\eji{0}{0}+3\eji{2}{0}+\eji{3}{0}$\newline
    ${}+\eji{1}{1}+3\eji{2}{1}+2\eji{3}{1}$\newline
    ${}+2\eji{0}{2}+\eji{2}{2}+\eji{3}{2})$
    \\[3pt]
8 & 0 & $\eji{0}{2}+\eji{2}{2}$ & 0
  & $2\eji{0}{0}+3\eji{2}{0}+\eji{3}{0}$\newline
    ${}+\eji{1}{1}+3\eji{2}{1}+2\eji{3}{1}$\newline
    ${}+3\eji{0}{2}+2\eji{2}{2}+\eji{3}{2}$
    & 0
  & $\tfrac 1 9 (2\eji{0}{0}+3\eji{2}{0}+\eji{3}{0}$\newline
    ${}+\eji{1}{1}+3\eji{2}{1}+2\eji{3}{1}$\newline
    ${}+3\eji{0}{2}+2\eji{2}{2}+\eji{3}{2})$
   \\
\midrule
&\multicolumn{2}{p{3.5cm}|}{$\sum_s \delta_{(s)}=20$}
  & \multicolumn{2}{p{3.5cm}|}{$\sum_s \Delta_{(s)}=5\ne 20$}
  & \multicolumn{2}{p{3.5cm}|}{$\sum_s \alpha_{(s)}=20$}
  \\
     \midrule
& \multicolumn{2}{p{3.5cm}|}{
   $\sum_s \delta_{(s)} V_{(s)}=$}
  &\multicolumn{2}{p{3.5cm}|}{$
  \sum_s \Delta_{(s)} L_{(s)}=$}
  &\multicolumn{2}{p{3.5cm}}{$
  \sum_s \alpha_{(s)} {\hat L}_{(s)}=$}
   \\
&
&
   $7\eji{0}{0}+11\eji{2}{0}+5\eji{3}{0}$\newline
    ${}+3\eji{1}{1}+7\eji{2}{1}+5\eji{3}{1}$\newline
    ${}+2\eji{0}{2}+\eji{2}{2}+\eji{3}{2}$
  &
  &
  $7\eji{0}{0}+11\eji{2}{0}+5\eji{3}{0}$\newline
    ${}+3\eji{1}{1}+7\eji{2}{1}+5\eji{3}{1}$\newline
    ${}+2\eji{0}{2}+\eji{2}{2}+\eji{3}{2}$
  &
  & $7\eji{0}{0}+11\eji{2}{0}+5\eji{3}{0}$\newline
    ${}+3\eji{1}{1}+7\eji{2}{1}+5\eji{3}{1}$\newline
    ${}+2\eji{0}{2}+\eji{2}{2}+\eji{3}{2}$
  \\
\bottomrule
\end{tabular}
\end{center}

%\noindent
%Note how $L_s$ grows by exactly $V_{(s)}$ at each step: newly introduced Ejis
%appear with coefficient~1; Ejis already present have their count incremented.
%$L_s$~with index~$s{+}1$ is thus the count of how many of the first $s{+}1$
%prefix elements contain each $\eji{j}{i}$.

%Calculate $\alpha_s = (s+1)(\delta_{(s)}-\delta_{(s+1)})$:
%\[
%  1,\; 2,\; 3,\; 0,\; 0,\; 6,\; 0,\; 8,\; 0.
%\]
\begin{tcolorbox}[notebox]
\textbf{Coincidence in this example.}  Every non-zero $\Delta$ (i.e.~every difference
between successive $\delta$ values)
happens to equal~$1$ here (e.g.\ $\delta_{(0)}-\delta_{(1)}=5-4=1$,
$\delta_{(5)}-\delta_{(6)}=2-1=1$, etc.), so the non-zero $\alpha_s$ values
equal $s{+}1$.  This is a coincidence of the chosen data, not a general rule:
$\alpha_s = (s{+}1)\times(\text{gap difference})$, and the gap difference can
be any non-negative integer.
\end{tcolorbox}

\begin{tcolorbox}[notebox]
\textbf{Verification: column sums and three-way equality.}
The integer columns $\delta_{(s)}$ and $\alpha_s$ have the same total:
\[
  \sum_{s=0}^{8}\delta_{(s)} \;=\; 5+4+3+2+2+2+1+1+0 \;=\; 20
  \;=\; 1+2+3+0+0+6+0+8+0 \;=\; \sum_{s=0}^{8}\alpha_s.
\]
The three weighted Eji sums are all equal and recover the
\texttt{balance} matrix of Step~B (the $L_{(s)}$ and ${\hat L}_{(s)}$ totals-row
entries show the value directly; $\sum_s\delta_{(s)}V_{(s)}$ is identical):
\begin{align*}
  \sum_s \delta_{(s)}\,V_{(s)}
  \;=\; \sum_s \Delta_{(s)}\,L_{(s)}
  \;=\; \sum_s \alpha_s\,{\hat L}_{(s)}
  &= 7\eji{0}{0}+11\eji{2}{0}+5\eji{3}{0} \\
  &\quad{}+3\eji{1}{1}+7\eji{2}{1}+5\eji{3}{1} \\
  &\quad{}+2\eji{0}{2}+\eji{2}{2}+\eji{3}{2},
\end{align*}
i.e.\ coefficient matrix
$\left(\begin{smallmatrix}7&0&11&5\\0&3&7&5\\2&0&1&1\end{smallmatrix}\right)$
(rows $=i=0,1,2$; columns $=j=0,1,2,3$), matching \texttt{balance} exactly.
\end{tcolorbox}

In many implementations the basis elements $L_{(s)}$ or ${\hat L}_{(s)}$ are stored as matrices.   For example, Eji\_LinComb
structures \url{https://github.com/kesterlester/Echinocoder/blob/main/Eji_LinComb.py} like those shown
below (only for each term without a zero coefficient) are used in
\url{https://github.com/kesterlester/Echinocoder/blob/main/C0HomDeg1_simplicialComplex_embedder_2_for_array_of_reals_as_multiset.py}.
In these the count matrix has rows representing
coordinates $i{=}0,1,2$; and columns representing vectors $j{=}0,1,2,3$.
\begin{center}
\small
\setlength{\tabcolsep}{4pt}
\renewcommand{\arraystretch}{1.2}
\begin{tabular}{cl l p{6.6cm}}
\toprule
$s$ & $\alpha_s$ & Eji\_LinComb structure
    &  ${\hat L}_{(s)}=L_{(s)}/(s{+}1)$ \\
\midrule
0 & 1
  & $\begin{pmatrix}1&0&1&1\\0&0&0&0\\0&0&0&0\end{pmatrix}$, idx 1
  & $\eji{0}{0}+\eji{2}{0}+\eji{3}{0}$ \\[8pt]
1 & 2
  & $\begin{pmatrix}1&0&2&1\\0&0&0&0\\0&0&0&0\end{pmatrix}$, idx 2
  & $\tfrac{1}{2}(\eji{0}{0}+2\eji{2}{0}+\eji{3}{0})$ \\[8pt]
2 & 3
  & $\begin{pmatrix}1&0&2&1\\0&1&1&1\\0&0&0&0\end{pmatrix}$, idx 3
  & $\tfrac{1}{3}(\eji{0}{0}+2\eji{2}{0}+\eji{3}{0}$\newline
    $\phantom{\tfrac{1}{3}(}+\eji{1}{1}+\eji{2}{1}+\eji{3}{1})$ \\[8pt]
5 & 6
  & $\begin{pmatrix}2&0&3&1\\0&1&3&2\\0&0&0&0\end{pmatrix}$, idx 6
  & $\tfrac{1}{6}(2\eji{0}{0}+3\eji{2}{0}+\eji{3}{0}$\newline
    $\phantom{\tfrac{1}{6}(}+\eji{1}{1}+3\eji{2}{1}+2\eji{3}{1})$ \\[8pt]
7 & 8
  & $\begin{pmatrix}2&0&3&1\\0&1&3&2\\2&0&1&1\end{pmatrix}$, idx 8
  & $\tfrac{1}{8}(2\eji{0}{0}+3\eji{2}{0}+\eji{3}{0}$\newline
    $\phantom{\tfrac{1}{8}(}+\eji{1}{1}+3\eji{2}{1}+2\eji{3}{1}$\newline
    $\phantom{\tfrac{1}{8}(}+2\eji{0}{2}+\eji{2}{2}+\eji{3}{2})$ \\
\bottomrule
\end{tabular}
\end{center}

\noindent
(Terms with $\alpha_s=0$ contribute nothing and are omitted.)

\subsection*{Step E: Canonicalise}

Sort the columns of each count matrix lexicographically (the column-sort
permutation ``forgets'' which vector was which, achieving multiset invariance).
After canonicalisation, $L_s$ becomes $\canon{L_s}$.

\subsection*{Step F: Hash and assemble}

Each $\canon{L_s}$ (along with its index) is hashed to a point $p_s \in
[0,1]^{19}$ (since $N_{\text{hash}} = 2{\times}9{+}1 = 19$).
The output is:
\[
  f(M) \;=\;
  \bigl[\;\underbrace{-3,\;2,\;1}_{3\text{ row min.}},\;
          1\,p_0 + 2\,p_1 + 3\,p_2 + 6\,p_5 + 8\,p_7\;\bigr]
  \;\in\; \R^{22}.
\]

% ═════════════════════════════════════════════════════════════════════════════
\section{Worked Example: Algorithm 1, $n=4$, $k=2$}
\label{sec:example-alg1}
% ═════════════════════════════════════════════════════════════════════════════

\[
  M \;=\; \left\{\;
    \begin{pmatrix}4\\2\end{pmatrix},\;
    \begin{pmatrix}-3\\5\end{pmatrix},\;
    \begin{pmatrix}8\\9\end{pmatrix},\;
    \begin{pmatrix}2\\7\end{pmatrix}
  \;\right\}
  \quad (n=4 \text{ vectors in } \R^2)
\]

\subsection*{Step B: Extract global minimum (and maximum)}

Global min $= -3$, global max $= 9$.  Both stored as anchors.  Balance has
all non-negative entries.

\subsection*{Steps C--D: Global sort and Abel summation}

The balance after subtracting the global minimum $-3$ from every entry is:
\[
  \left\{\;\begin{pmatrix}7\\5\end{pmatrix},\;
            \begin{pmatrix}0\\8\end{pmatrix},\;
            \begin{pmatrix}11\\12\end{pmatrix},\;
            \begin{pmatrix}5\\10\end{pmatrix}\;\right\}
\]
Sorted globally descending (value: label):
\[
  12(\eji{2}{1}),\;11(\eji{2}{0}),\;10(\eji{3}{1}),\;8(\eji{1}{1}),\;
  7(\eji{0}{0}),\;5(\eji{0}{1}),\;5(\eji{3}{0}),\;0(\eji{1}{0}).
\]
The 7 gaps are $1,1,2,1,2,0,5$.  After global re-sort by gap descending and
second Abel summation: $\alpha_0=3,\;\alpha_2=3,\;\alpha_5=6$ (others zero).

$N_{\text{hash}} = 2(8{-}1){+}1 = 15$. Output:
\[
  f(M) \;=\;
  \bigl[\;3\,p_0 + 3\,p_2 + 6\,p_5,\;9,\;{-3}\;\bigr]
  \;\in\; \R^{17},
\]
numerically (from the code's debug output):
\[
  [\,7.07,\;6.91,\;2.78,\;3.50,\;1.01,\;5.51,\;4.07,\;4.90,\;10.02,\;
   9.51,\;5.13,\;2.97,\;8.04,\;5.34,\;7.41,\;9,\;{-3}\,].
\]

% ═════════════════════════════════════════════════════════════════════════════
\section{Output Size and the Whitney Embedding Theorem}
\label{sec:whitney}
% ═════════════════════════════════════════════════════════════════════════════

Both algorithms produce output of size $\approx 2nk$, which is not overhead to
be engineered away — it is intrinsic.  The input space $\SP{n}{k}$ is an
$nk$-dimensional manifold.  The \textbf{Whitney Embedding Theorem} states that
any smooth $d$-dimensional manifold embeds smoothly into $\R^{2d}$, and $2d$
is tight in general.  With output $\approx 2nk$, both algorithms sit right at
the Whitney bound.

The factor of $\approx 2$ is in fact the price of the cone-dimension mechanism
in Step~F: encoding $m$ positive coefficients together with the identities of
$m$ basis elements inside a single point requires $\R^{2m+1}$.  This same
mechanism is what delivers $C^0$ continuity (shared faces at zero-coefficient
transitions).  The factor of two is the minimum price of a smooth
embedding, and the algorithms pay almost exactly that.

Whether $2nk$ (or $2nk+1$, or $2nk+1-k$) is optimal for the specific
manifold $\SP{n}{k}$, or whether the known structure of that space admits a
tighter construction, remains open.

\end{document}
